\section{Augmentation: o build que roda o framework}
\label{sec:augmentation}

No Quarkus, o build não só compila: ele executa a parte pesada do
framework. Essa etapa se chama \termo{augmentation}{\enquote{Aumento}: a
fase do build em que as extensões do Quarkus analisam a aplicação e
geram o código que será executado em runtime. Na documentação aparece
também como \emph{build-time augmentation}.}, e o guia
\citetitle{doc:quarkus-writing-extensions} dá a regra que orienta todas
as extensões: \enquote{As much as possible favor doing work at build
time}. O motivo vem logo depois: quanto mais trabalho for feito ali,
menores ficam as aplicações e mais rápido elas sobem. A introdução do
mesmo guia acrescenta que os metadados são processados uma só vez, no
build.

\subsection*{Extensões em duas metades}

Toda extensão tem dois módulos. O módulo \emph{deployment} roda só no
build; o módulo \emph{runtime} vai para o artefato final. O guia é
taxativo: o módulo runtime nunca pode depender de um artefato de
deployment. É por isso que o parser de anotações, a validação do grafo
de beans e o gerador de código não pesam no seu container: eles ficaram
na máquina de build.

\subsection*{Build steps e build items}

O trabalho do deployment é dividido em \termo{build steps}{Passos do
build: métodos anotados com \texttt{@BuildStep} numa classe de
processamento da extensão.}, que se comunicam por \termo{build
items}{Objetos que um passo produz e outros consomem; são a única forma
de as extensões trocarem informação no build.}. Dois detalhes do guia
importam para quem só usa o Quarkus:
\begin{itemize}
\item Um build step só roda se produzir algo de que outro passo precisa,
  direta ou transitivamente. Extensão que você não usa não custa nada no
  build nem no runtime.
\item Os processadores de build \enquote{should not attempt to load any
  application classes}: eles leem a aplicação pelo índice
  (seção~\ref{sec:jandex}), não carregando as classes. Carregar uma classe
  dispararia linking e, no primeiro uso ativo, o \texttt{<clinit>}
  (seção~\ref{sec:class-loading}).
\end{itemize}

\subsection*{Recorders: STATIC\_INIT e RUNTIME\_INIT}

Parte do trabalho não pode ser feito no build (abrir um pool de conexões,
por exemplo), mas pode ser \emph{preparado} nele. Para isso existem os
\termo{recorders}{Gravadores: classes do módulo runtime cujas chamadas,
feitas durante o build, são gravadas e reproduzidas na subida da
aplicação.}. No build, a extensão recebe um proxy do recorder que grava
as chamadas; segundo o guia, essa gravação é usada para gerar bytecode que
faz a mesma sequência de chamadas em runtime. O momento da reprodução
depende da anotação:

\begin{description}
\item[\texttt{@Record(STATIC\_INIT)}] O bytecode roda num método de
  inicialização estática da classe principal gerada. Nesse momento é
  proibido ler configuração de runtime.
\item[\texttt{@Record(RUNTIME\_INIT)}] O bytecode roda a partir do
  \texttt{main} da aplicação, e aí a configuração de runtime já está
  disponível.
\end{description}

A divisão faz sentido à luz do executável nativo. O guia
\citetitle{doc:graalvm-native-basics} diz que, se uma classe é
inicializada em build time, os campos \texttt{static} dela são salvos no
binário, e o primeiro uso em runtime não dispara inicialização. Tudo o
que roda em STATIC\_INIT pode, portanto, virar parte do executável já
pronto; ler uma variável de ambiente ali significaria congelar o valor
da máquina de build. O capítulo~\ref{cap:decisoes-quarkus} volta a isso.

\subsection*{Vendo o build trabalhar}

O guia lista duas opções úteis para curiosos:
\texttt{-Dquarkus.debug.print-startup-times=true}, passado ao
\emph{iniciar} a aplicação, imprime quanto cada tarefa de startup gravada
por um recorder leva em runtime, e
\texttt{-Dquarkus.builder.graph-output=build.dot}, passado no build, grava o
grafo de dependências entre os passos. Rodar as duas uma vez num projeto
seu é a forma mais rápida de entender o que cada extensão faz.
